\chapter{Jenkins: The CI Server}
\label{ch:jenkins}

\section{Architecture}

Jenkins is a Java web application --- the \emph{controller} --- that
stores jobs, credentials, plugins and build history in one directory,
\code{JENKINS\_HOME}, and executes builds either on itself (via
\emph{executors}, its parallel-build slots) or on remote \emph{agents}.
The enterprise pattern is a controller that builds nothing, farming
work out to disposable agents (often pods, via the Kubernetes plugin).

\begin{decision}[builds on the controller, Jenkins outside the cluster]
Two homelab inversions of enterprise practice, both CPU arithmetic.
Agent pods would cost scheduling and image pulls per build on the same
two cores the build needs; the controller's two executors spend those
cycles on Maven instead. And Jenkins runs in Docker on the host, not in
k3s: if you break the cluster --- and while learning, you will ---
the thing that can rebuild the cluster must not have died with it.
Both revert cleanly at enterprise scale: in-cluster Jenkins with pod
agents is Chapter~26 material, and the migration itself is the lesson.
\end{decision}

\section{The controller image, annotated}

Rather than clicking tools into a stock Jenkins, the project bakes its
toolchain into an image --- \code{deploy/jenkins/Dockerfile}:

\begin{lstlisting}[language=dockerfile]
FROM jenkins/jenkins:lts-jdk21          (*@\co{1}@*)
USER root
RUN ... temurin-25-jdk ...              (*@\co{2}@*)
RUN ... apache-maven-3.9.9 ...          (*@\co{3}@*)
RUN wget -qO /usr/local/bin/kubectl ... (*@\co{4}@*)
ENV JDK25_HOME=/usr/lib/jvm/temurin-25-jdk-amd64  (*@\co{5}@*)
USER jenkins
RUN jenkins-plugin-cli --plugins \
      git workflow-aggregator ... credentials-binding  (*@\co{6}@*)
\end{lstlisting}

\begin{description}
  \item[\co{1}] LTS Jenkins running on JDK\,21 --- the version
  \emph{Jenkins itself} runs on.
  \item[\co{2}] The project targets Java\,25, so a second JDK is
  installed for \emph{builds}. Note \code{JAVA\_HOME} is \emph{not}
  changed globally \co{5} --- the Jenkinsfile sets it per pipeline.
  Controller runtime and build toolchain are independent concerns; this
  Dockerfile keeps them visibly separate.
  \item[\co{3}] Maven pinned to 3.9.9 --- distro packages lag.
  \item[\co{4}] \code{kubectl}: the deploy stage is just this binary
  plus a kubeconfig.
  \item[\co{6}] Plugins pinned in the image, not clicked in the UI:
  pipeline support, git, JUnit reporting, and
  \code{credentials-binding} (used by \code{withCredentials} in
  Chapter~23). The whole point: \emph{the CI server itself is now
  reproducible} --- destroy the container, rebuild from the Dockerfile,
  and only \code{JENKINS\_HOME} (a named volume) carries state.
\end{description}

One run-time flag from \code{run.sh} deserves its own paragraph:
\code{--network host}. Jib inside Jenkins pushes to
\code{localhost:5000} --- and inside a bridge-networked container,
\code{localhost} would be the \emph{container}. Host networking makes
Jenkins share the host's network namespace, so its
\code{localhost:5000} is the same registry k3s pulls from
(Chapter~13's naming contract). Not a shortcut --- a requirement of the
image-naming scheme.

\section{Credentials}

Jenkins stores secrets in an encrypted store, referenced by ID;
pipelines bind them to environment variables only for the duration of a
block (\code{withCredentials}), and the log masker redacts their
values from console output. This project stores one credential:
\code{kubeconfig-ts}, a file credential holding the kubeconfig Jenkins
uses to deploy --- which is the interesting part.

\section{The deploy identity: a ServiceAccount, not admin}

The lazy path is giving Jenkins the k3s admin kubeconfig ---
\code{cluster-admin}, licensed to delete every namespace. Instead,
\code{deploy/k8s/jenkins-rbac.yaml} builds a minimal identity:

\begin{lstlisting}[language=yaml]
kind: ServiceAccount          # the identity
metadata: { name: jenkins, namespace: ts }
---
kind: Role                    # namespaced permissions
metadata: { name: jenkins-deployer, namespace: ts }
rules:
  - apiGroups: ["apps"]
    resources: ["deployments", "statefulsets"]
    verbs: ["get", "list", "watch", "create", "patch", "update"] (*@\co{1}@*)
  - apiGroups: ["apps"]
    resources: ["replicasets"]
    verbs: ["get", "list", "watch"]                              (*@\co{2}@*)
  # + services, PVCs, secrets, configmaps, ingresses, pods/log ...
  - apiGroups: ["kafka.strimzi.io"]                              (*@\co{3}@*)
    resources: ["kafkas", "kafkanodepools", "kafkatopics"]
    verbs: ["get", "list", "create", "patch", "update"]
---
kind: RoleBinding             # identity <- permissions
\end{lstlisting}

\begin{description}
  \item[\co{1}] The verbs are derived from the pipeline's actual
  commands: \code{kubectl apply} needs get+create+patch;
  \code{set image} needs get+patch. \emph{No delete anywhere}: CI rolls
  forward; a human deletes. A compromised pipeline can thus corrupt
  deployments in \code{ts} but cannot wipe them, and cannot see beyond
  the namespace --- a Role, not a ClusterRole, is the blast-radius
  decision.
  \item[\co{2}] Discovered the honest way: \code{rollout status}
  watches the Deployment's ReplicaSets, so read access to them is
  required --- the kind of permission you find by running the pipeline
  against the Role and reading the first \code{Forbidden}.
  \item[\co{3}] Added with milestone~6: the pipeline may declare Kafka
  clusters and topics (namespaced custom resources), but installing
  the operator and its CRDs --- cluster-scoped --- stays an admin act
  (Chapter~25). RBAC grows one rule per API group the pipeline touches;
  the Role file is a changelog of what CI has been trusted with.
\end{description}

The kubeconfig stored in Jenkins simply pairs the cluster address with
this ServiceAccount's token instead of the admin certificate.

\begin{pitfall}[Forbidden: cannot create resource ``namespaces'']
Symptom: the deploy stage dies with a Forbidden error on
\code{namespaces} although the RBAC file looks complete. Cause:
namespaces are \emph{cluster-scoped}, and a namespaced Role cannot
grant them --- which is exactly why \code{namespace.yaml} is excluded
from the kustomization (Chapter~20) and the namespace is created once
by an admin. RBAC debugging in general:
\code{kubectl auth can-i patch deployments -n ts --as
system:serviceaccount:ts:jenkins} answers ``would this be allowed''
without running the pipeline.
\end{pitfall}

\begin{handson}[prove the boundary exists]
\begin{lstlisting}[language=bash]
kubectl auth can-i patch deployments -n ts \
  --as system:serviceaccount:ts:jenkins        # yes
kubectl auth can-i delete deployments -n ts \
  --as system:serviceaccount:ts:jenkins        # no
kubectl auth can-i list pods -n kube-system \
  --as system:serviceaccount:ts:jenkins        # no
\end{lstlisting}
\end{handson}

\begin{quiz}
\begin{enumerate}
  \item What lives in \code{JENKINS\_HOME}, and why does baking plugins
  into the image plus a volume for home make the CI server disposable?
  \item Why does the controller run two JDKs, and where is each
  selected?
  \item Reconstruct the RBAC verb list for a pipeline that only ever
  runs \code{kubectl apply -k} and \code{rollout status} --- which verb
  in the project's Role could then be dropped?
  \item Explain \code{--network host} on the Jenkins container in terms
  of Chapter~13's image-naming contract.
  \item ServiceAccount, Role, RoleBinding: which is identity, which is
  capability, which is the join --- and why is Role-not-ClusterRole the
  security headline?
\end{enumerate}
\end{quiz}
